% NOTE: the official solution refers to "figure 1" (maximum-brightness
% configuration) and "Figure 2" (total eclipse, projected circle of radius b),
% but no such figures are printed in the source solutions PDF
% (2009_th_sol.pdf, pp. 54--55).

Ignoring temperature variation on the stars' surface, the brightness of the
star system will be proportional to the projected surface of both stars on the
plane of the sky. Maximum brightness will occur when the two stars are seen
like figure 1 and minimum light will happen when one of the stars is in total
eclipse and the projected surface of the other is a circle with radius $b$
(Figure 2). In maximum brightness
\[ I_{max} \propto 2\pi a b \]
\hfill(3 points)

In minimum brightness
\[ I_{min} \propto \pi b^2 \]
\hfill(3 points)

So
\[ \Delta m = -2.5\log\frac{I_{max}}{I_{min}}
   = -2.5\log\frac{2\pi ab}{\pi b^2}
   = -2.5\log 4 \]
\hfill(2 points)
\[ \Delta m = -1.5 \]
\hfill(2 points)
